\documentclass{article}
\usepackage[utf8]{inputenc}
\usepackage[T1]{fontenc}
\usepackage{geometry}
\usepackage{enumitem}
\usepackage{hyperref}
\usepackage{xcolor}
\usepackage{titlesec}
\usepackage{times}

\geometry{a4paper, margin=1in}

\title{\textbf{The Hidden Riches of Punjab}}
\subtitle{Culture, Manuscripts, and Linguistic Heritage}
\author{}
\date{}

\titleformat{\section}
  {\normalfont\Large\bfseries}{\thesection}{1em}{}
\titleformat{\subsection}
  {\normalfont\large\bfseries}{\thesubsection}{1em}{}

\begin{document}

\maketitle

\begin{center}
\textit{An exploration of traditions, textual treasures, and linguistic diversity beyond mainstream narratives.}
\end{center}

\section{Vibrant Culture}
The culture of Punjab is a dynamic tapestry of folk music, dance, oral traditions, festivals, and a resilient community spirit. While Punjabi culture is often reduced to \textit{bhangra} and \textit{naan} in popular imagination, its deeper layers remain less visible to the broader Indian population. At its heart lies the \textbf{folk heritage}: \textit{Jugni} (spiritual folk ballads), \textit{Heer} (epic romances like \textit{Heer Ranjha} by Waris Shah), and the \textit{dhadhi} \textit{vaar} (ballad-singers narrating heroic tales). Festivals such as \textit{Lohri}, \textit{Maghi}, and \textit{Teej} are not merely seasonal events but are interwoven with agrarian cycles and collective memory.

The \textbf{Sikh Gurdwara} system, while religious, also serves as a cultural nucleus promoting \textit{seva} (selfless service) and \textit{langar} (community kitchen)—a tradition that embodies equality. Moreover, Punjab's folk instruments—the \textit{tumbi}, \textit{chimta}, \textit{algoza}, and \textit{sarangi}—carry melodic structures that predate classical music codification. The \textbf{phulkari} embroidery tradition, with its geometric motifs and deep symbolism, represents not just craft but a language of familial bonds and regional identity.

\begin{quote}
\textit{``The soul of Punjab resides not in its cities but in its folk lores, where every village has its own \textit{dhadhi} and every soil sings a different \textit{boli}.''} \cite{singh2018folk}
\end{quote}

\section{Manuscripts}
Punjab’s manuscript tradition spans multiple scripts, religions, and centuries, yet remains largely inaccessible to the public due to preservation issues and linguistic barriers. The most celebrated repository is the \textbf{Guru Granth Sahib}, which contains hymns in over a dozen linguistic registers, recorded in the \textbf{Gurmukhi} script. However, a vast corpus of pre-colonial manuscripts exists in \textbf{Takri}, \textbf{Landey}, and \textbf{Perso-Arabic} scripts.

Key collections include:
\begin{itemize}
    \item \textbf{Janamsakhis}: Hagiographical manuscripts narrating the life of Guru Nanak, often illuminated with miniature paintings. The \textit{Bhai Bala} and \textit{Puratan} versions exist in unique proto-Gurmukhi scripts.
    \item \textbf{Pothis and Gutkas}: Handwritten compilations of hymns, some dating back to the 17th century, held in private \textit{deras} and the Punjab State Archives.
    \item \textbf{Sufi literature}: Manuscripts of Shah Hussain, Bulleh Shah, and Sultan Bahu, primarily in the Persian and Shahmukhi scripts, preserved in shrines like the \textit{Shrine of Hazrat Data Ganj Bakhsh} in Lahore and in private collections in East Punjab.
    \item \textbf{Heroic poetry}: \textit{Vaars} (ballads) composed by poets like \textit{Shah Mohammad} during the Anglo-Sikh wars, surviving in fragile handwritten copies.
\end{itemize}

In Indian Punjab, institutions such as the \textbf{Punjab Digital Library} (PunjabDigitalLibrary.com) and the \textbf{Sikh Reference Library} (post-1984 reconstruction efforts) are working to digitize over 5 million folios. Yet, the majority of these manuscripts remain in village \textit{deras} (shrines) and family collections, unclassified and at risk.

\section{Local Languages}
Punjab’s linguistic landscape is far more diverse than the singular label “Punjabi” suggests. The region is home to a continuum of languages and dialects that vary by region, caste, and religion. These include:

\begin{itemize}
    \item \textbf{Majhi}: The standard dialect, spoken in the central districts (Amritsar, Lahore, etc.), considered the prestige variety.
    \item \textbf{Doabi}: Spoken in the Doaba region between the Beas and Sutlej rivers, with distinct phonetic and lexical features.
    \item \textbf{Malwai}: Prevalent in the Malwa region (southern Punjab), known for its robust vocabulary and unique proverbs.
    \item \textbf{Pothohari} and \textbf{Hindko}: Spoken in the northern highlands and western districts, often treated as separate languages but sharing deep historical roots with Punjabi.
    \item \textbf{Puadhi}: A transitional dialect between Punjabi and Haryanvi, spoken in the southeastern districts.
    \item \textbf{Shahmukhi}: The Perso-Arabic script used by Punjabi Muslims across the border, which, along with Gurmukhi, creates a digraphic culture often invisible to monolingual speakers.
\end{itemize}

Additionally, \textbf{Sansi}, \textbf{Bazigar}, and \textbf{Lubanki} — languages of nomadic and denotified tribes — survive as marginalized linguistic communities. The \textbf{2011 Census of India} records Punjabi as the 11th most spoken language nationally, but it aggregates all dialects under a single head, erasing local linguistic diversity. Scholars such as Dr. Harjit Singh Gill have documented these dialectal variations, emphasizing that “the vitality of Punjab lies in its \textit{boliyan} (folk couplets) and their subtle regional inflections” \cite{gill2010dialect}.

\subsection*{Preservation and Challenges}
Despite a vibrant oral and textual culture, institutional neglect, the partition of 1947 (which divided archival holdings), and the dominance of Hindi and English in education have rendered much of this heritage hidden. Recent grassroots initiatives, including community-led archives and digital humanities projects, aim to recover these cultural layers.

\begin{thebibliography}{9}
\bibitem{singh2018folk}
Singh, K. (2018). \textit{Folk Traditions of Punjab: Beyond Bhangra}. Patiala: Punjabi University Press.

\bibitem{gill2010dialect}
Gill, H. S. (2010). “The Dialectal Structure of Punjabi.” In \textit{Journal of Punjabi Studies}, 17(1), 45–62.

\bibitem{punjabdigital}
Punjab Digital Library. (n.d.). \textit{Manuscripts and Heritage}. Retrieved from \url{https://www.panjabdigilib.org}

\bibitem{grewal2004sikh}
Grewal, J. S. (2004). \textit{Historical Writings on the Sikhs (1784–2011)}. New Delhi: Manohar Publishers. (Chapter on manuscript traditions).

\bibitem{shackle2003punjabi}
Shackle, C. (2003). “Punjabi.” In \textit{The Indo-Aryan Languages}, edited by G. Cardona and D. Jain. London: Routledge. (Provides comprehensive overview of dialects and scripts).

\end{thebibliography}

\end{document}